% One event through the system, and where latency is defined.
\begin{tikzpicture}[font=\sffamily\scriptsize,
  mlab/.style={fill=white, inner sep=1.5pt, font=\sffamily\scriptsize}]
\foreach \x/\name/\id in {0/timer ISR/isr, 3.2/scheduler/sch, 6.4/sensor task/sen,
                          9.6/feature task/fea, 12.8/link task/lnk}{
  \node[actor, text width=20mm] (\id) at (\x,0) {\name};
  \draw[lifeline] (\id.south) -- ++(0,-8.4);
}

\draw[msg] (-1.4,-1.0) -- node[mlab, above]{release edge} (0,-1.0);
\node[activation, minimum height=4mm] at (0,-1.3) {};
\draw[msg] (0,-1.5) -- ++(0.9,0) -- ++(0,-0.4)
  node[mlab, right, xshift=1mm]{timestamp, then notify} -- ++(-0.9,0);
\draw[msg] (0,-2.5) -- node[mlab, above]{yield from ISR} (3.2,-2.5);
\node[activation, minimum height=4mm] at (3.2,-2.8) {};
\draw[msg] (3.2,-3.3) -- node[mlab, above]{context switch} (6.4,-3.3);
\node[activation, minimum height=6mm] at (6.4,-3.6) {};
\draw[msg] (6.4,-3.9) -- ++(0.9,0) -- ++(0,-0.4)
  node[mlab, right, xshift=1mm]{timestamp, then read the FIFO} -- ++(-0.9,0);
\draw[msg] (6.4,-5.2) -- node[mlab, above]{bytes into the stream buffer} (9.6,-5.2);
\node[activation, minimum height=5mm] at (9.6,-5.5) {};
\draw[msg] (9.6,-6.4) -- node[mlab, above]{payload into the queue} (12.8,-6.4);
\node[activation, minimum height=4mm] at (12.8,-6.7) {};
\draw[reply] (12.8,-7.6) -- node[mlab, below]{result line, from a task} (0,-7.6);

% the measured interval
\draw[tspan] (-1.1,-1.0) -- (-1.1,-3.9);
\node[mlab, rotate=90, anchor=south] at (-1.25,-2.45) {the measured interval};

\node[umlnote, text width=54mm, anchor=north west] at (-1.6,-9.2)
  {Latency is defined from the release edge to the first instruction of the
   task body, so it includes the handler, the scheduler and the context switch.
   Measuring only to the end of the handler reports a smaller number that no
   deadline depends on.};
\node[umlnote, text width=54mm, anchor=north west] at (7.2,-9.2)
  {The result line is printed from a task and never from inside the interval
   being measured. Printing from the handler is the fourth most likely cause of
   a tail that looks wrong.};
\end{tikzpicture}
